\clearpage
\fancyhead[L]{\small BODE. ASTRONOMISCHES JAHRBUCH.}
\hypertarget{native-624-lower}{}
\begin{center}\small Göttingische gelehrte Anzeigen. Stück 23. Seiten 225..232. 7. Februar 1820.\\[3pt]\rule{.72\linewidth}{.3pt}\end{center}

\emph{Astronomisches Jahrbuch für das Jahr {\upshape 1822}, nebst einer Sammlung der neuesten Abhandlungen, Beobachtungen und Nachrichten aus dem Gebiet der astronomischen Wissenschaften. Berechnet und herausgegeben von} Dr. J. E. \textsc{Bode}, \emph{königlicher Astronom usw.} Berlin. Beim Verfasser und in Kommission bei F. Dümmler. 260 Seiten im Oktavformat, nebst einer Kupfertafel.

\begingroup\rightskip=0pt plus 1em
Von den zwei Sonnen- und zwei Mondfinsternissen, die sich im Jahr 1822 ereignen, ist in Europa nur eine der Mondfinsternisse sichtbar. Auch der Merkurtransit am 5. November kann in Europa nicht beobachtet werden. Die Einrichtung des Jahrbuchs ist gegenüber den früheren Jahrgängen gleich geblieben. Nur die Verfinsterungen der Jupitermonde sind zur Erleichterung der Arbeit nach den Tafeln von \textsc{Wargentin} berechnet worden, und nur bei den ersten Monden wurden gelegentlich die Tafeln von \textsc{Delambre} mit herangezogen.\par\endgroup

\begingroup\brokenpenalty=10000\relax
Den Anfang unter den Abhandlungen im Anhang macht ein Versuch von Prof. \textsc{Buttmann}, die astrognostischen Namen einzelner Sterne festzulegen. Dieser gelehrte Philologe \SourceLeaf{625}stellt hier ein Verzeichnis von 80 Sternnamen auf, die größtenteils der arabischen Sprache entlehnt sind. Einige waren schon gebräuchlich oder doch nicht unbekannt, andere berichtigt oder an eine andere Stelle gesetzt, und wieder andere sind ganz neu gebildet (z. B. Clavus für $\beta$ im Stier, Meion für $\varkappa$ im Orion, 
Sectatrix für $\beta$ im Fuhrmann). Unter den Astronomen ist eigentlich nur eine sehr kleine Zahl von Eigennamen für die Fixsterne \emph{wirklich gebräuchlich} (etwa ein Dutzend), und die Versuche einiger, ältere arabische Namen wieder einzuführen, haben keinen Anklang gefunden. Wollte man dennoch einmal eine größere Zahl in Gebrauch bringen, wäre es wohl am zweckmäßigsten, sich einfach und \emph{ohne jede Änderung} an die Namen zu halten, die \textsc{Piazzi} in seinen neuen Sternkatalog aufgenommen hat. Jede willkürliche Änderung könnte, statt eine Verbesserung zu bringen, nur zu einer Verwirrung in der Benennung führen. --- Verfinsterungen der Jupitermonde, Sternbedeckungen, Mondfinsternisse und Oppositionen von Vesta, Uranus, Jupiter, Saturn und Pallas, im Jahr 1818 beobachtet von Prof. \textsc{Bürg}. --- Bestimmung von Algols Lichtperiode aus neueren Beobachtungen und Berechnung seines Lichtminimums für 1820 bis 1822 durch Prof. \textsc{Wurm}. Diese Untersuchung zeigt, dass sich die Länge der Periode seit der Entdeckung der Veränderlichkeit des Sterns überhaupt nicht merklich geändert hat. --- Astronomische Beobachtungen an der Prager Sternwarte von den Herren \textsc{David} und \textsc{Bittner}. --- Ebenfalls von Prof. \textsc{Bürg} Beobachtungen der Juno und Sternbedeckungen im Jahr 1819, nebst astronomischen Nachrichten. Die hier vorliegende Erwiderung auf einen Angriff auf einen Zeitungsartikel über einen der Kometen von 1818 oder eine falsche Auslegung dieses Artikels (der Angriff oder die falsche Auslegung findet sich in der in Genua erscheinenden Correspondance astronomique) weckt den Wunsch, dass persönliche Bemerkungen aus astronomischen Erörterungen stets ferngehalten werden mögen. --- In Prag beobachtete Sternbedeckungen und Verfinsterungen der Jupitermonde von Prof. \textsc{Hallaschka}. --- Ältere Beobachtungen des Uranus von \textsc{Lemonnier} (siehe Jahrgang 1819, S. 1909). --- Über die Sonnenfinsternis vom 7. September 1820 für Deutschland und die angrenzenden Länder von Prof. \textsc{Littrow}. Hier werden Anfang und Ende für 230 Orte in Deutschland und den benachbarten Ländern angegeben, ebenso die durch eine Anzahl von Punkten bestimmten nördlichen und südlichen Grenzlinien der Zone, innerhalb derer die Finsternis ringförmig erscheinen wird. Der größere Teil Deutschlands liegt innerhalb dieser Zone. Die südliche Grenze verläuft nahe bei Aachen, Straßburg, Freiburg und Zürich, die nördliche nahe bei Schwerin, Wittenberg und Tabor in Böhmen. Es wird sehr wünschenswert sein, dass die Dauer der ringförmigen Finsternis an recht vielen Orten, besonders an solchen in der Nähe der Grenzen, beobachtet wird, .wozu lediglich ein gutes Fernrohr und eine Sekundenuhr erforderlich sind. --- Beobachtungen der Vesta, der Pallas, des Uranus, des Saturn und des Jupiter sowie der Sonnenfinsternis im Jahr 1818, in Wilna von Prof. \textsc{Sniadecky}. --- Über die geografische Länge der Berliner Sternwarte. Diese Länge wurde früher stets zu denen mit stark schwankenden Werten gezählt. Seit die Sternwarte jedoch mit besseren Hilfsmitteln zur Zeitbestimmung ausgestattet ist, zeigen auch die Beobachtungen eine befriedigendere Übereinstimmung. Aus 23 Sternbedeckungen und 4 Sonnenfinsternissen, die von 1802 bis 1815 beobachtet wurden, folgt im Mittel die Länge $44^{\mathrm{m}}10^{\mathrm{s}}5$ bezogen auf Paris. --- Beobachtete Oppositionen des Mars, des Uranus und des Jupiter 1817, 1818, ferner Beobachtungen der Sonnenfinsternis von 1818 und des \textsc{Enke}'schen Kometen, von Herrn \textsc{Derflinger} in Kremsmünster. --- Beobachtung einer Bedeckung des Mars durch den Mond, in Wien von Prof. \textsc{Bürg}. --- Astronomische Beobachtungen an der Berliner Sternwarte im Jahr 1818. Einer Berichtigung bedarf die hier aufgestellte übereilte Behauptung, dass die Bestimmung der heliozentrischen Position aus einer nahe der Opposition beobachteten geozentrischen Position wegen der kleinen Winkel des Dreiecks unsicher werde und dass es sogar vorteilhafter sei, Beobachtungen zugrunde zu legen, die weiter von der Opposition entfernt sind. --- Beobachtungen eines Kometen 1818 und 1819 von Prof. \textsc{Bessel}, ferner Planetenoppositionen, eine Sonnenfinsternis und Sonnenwenden. --- Von Dr. \textsc{Olbers} über die höchst bemerkenswerte \SourceLeaf{626} von Prof. \textsc{Enke} gemachte und in unseren Blättern (St. 28 u. 83 v. v. J. [S. 417 u.f. d. B.]) zuerst bekannt gegebene Entdeckung, dass einer der Kometen von 1805 mit einem der Kometen von 1818 identisch war. Dr. \textsc{Olbers} fügte dieser großen Entdeckung die ebenso wichtige Bemerkung hinzu, dass auch der Komet von 1795 höchstwahrscheinlich derselbe Komet war und dass auch die beiden in der Connaissance des tems 1819 bekannt gegebenen Kometenbeobachtungen von 1786 eben diesen Kometen betreffen. Ferner teilt dieser ausgezeichnete Astronom hier seine ersten Beobachtungen des im Juli dieses Jahres erschienenen Kometen mit. --- Der folgende schöne Aufsatz von Prof. \textsc{Enke} gibt nun einen ausführlichen Bericht über seine Arbeit zu jenem Kometen mit kurzer Umlaufzeit und erhebt die oben genannte wahrscheinliche Vermutung von Dr. \textsc{Olbers} zur moralischen Gewissheit. Mitte Mai 1822 wird dieser Komet nun erneut sein Perihel erreichen. Unglücklicherweise wird er zuvor wegen seiner Stellung zur Erde äußerst lichtschwach sein, sodass es schwierig sein wird, ihn bis dahin zu beobachten. Nach dem Durchgang durch das Perihel hingegen wird er zwar hell genug sein, aber zu weit südlich stehen, um in Europa sichtbar zu werden. Dann wird er sich jedoch auf der südlichen Hälfte der Erdkugel leicht beobachten lassen, sofern es dort nur nicht an Beobachtern fehlt. Bei Gelegenheiten wie dieser wird ebenso wie bei so vielen anderen das Bedürfnis nach einer \emph{südlichen} Sternwarte deutlich spürbar. Wir fühlen uns aus diesem Anlass gedrängt, unsere tiefe Überzeugung auszusprechen, dass man nach der großen Unterstützung, die die Astronomie in jüngster Zeit in Europa erhalten hat, dieser Wissenschaft keinen größeren Dienst mehr leisten könnte als durch die Errichtung einer gut ausgestatteten und mit einem einsichtsvollen, erfahrenen und tatkräftigen Astronomen zu besetzenden Sternwarte in der südlichen Hemisphäre, etwa am Kap der Guten Hoffnung. --- Von demselben Astronomen stammen Beobachtungen des Kometen von 1819 und die Elemente seiner Bahn, ferner Beobachtungen des Uranus, des Saturn, des Jupiter, der Ceres und der Juno 1818 und 1819. --- Lösung der Aufgabe, die Deklination eines Gestirns ohne Winkelmessinstrument allein mit einem Fernrohr zu bestimmen, von Prof. \textsc{Fischer} (aus der Dauer eines Durchgangs durch das Gesichtsfeld, bei dem Eintritt und Austritt an den Endpunkten eines Durchmessers erfolgen). --- Geometrischer Lauf der Ceres 1819 und 1820, berechnet von Dr. \textsc{Westphal}, und geozentrischer Lauf der Juno 1820, berechnet von Prof. \textsc{Nicolai}. --- Beobachtungen des Kometen vom Juli 1819 an der Berliner Sternwarte (sie reichen vom 2. bis 27. Juli). --- Von Prof. \textsc{Nicolai}, Beobachtungen der Juno an der Mannheimer Sternwarte und neue Bahnelemente, die aus diesen und aus Göttinger Beobachtungen abgeleitet sind. Ferner die ersten Beobachtungen des Polarsterns mit dem wiederhergestellten dreifüßigen Multiplikationskreis von \textsc{Reichenbach}, die in sehr guter Übereinstimmung die Polhöhe der Mannheimer Sternwarte mit $49^\circ\,29'\,12''95$ ergeben. Schließlich Beobachtungen und Elemente des Kometen von 1819. --- Die darauf folgenden Bemerkungen eines ungenannten Verfassers über den Gebrauch der Libelle und des Lots zur Justierung astronomischer Instrumente enthalten nichts Neues und beweisen nur, dass es darauf ankommt, wie Libellen an einem Instrument angebracht sind, dass sie eine ihrem Zweck angemessene Empfindlichkeit haben müssen und dass sie einen vorsichtigen Beobachter erfordern. Der hier gemachte Vorschlag, mit einem in einem Brunnen aufgehängten Lot zu prüfen, ob die Richtung der Schwerkraft veränderlich oder unveränderlich ist, dürfte kaum zu etwas anderem führen, als die \emph{Unerreichbarkeit absoluter Festigkeit} bei jeder Aufstellung von Messapparaten zu beweisen. --- Auszug aus einem Schreiben von Herrn General von \textsc{Lindener} aus Glatz. Da bei dem im Juli dieses Jahres erschienenen Kometen durch einen glücklichen Zufall die untere Konjunktion am Vormittag des 26. Juni mit dem Durchgang durch den aufsteigenden Knoten zusammenfiel, war es sehr erwünscht, dass dieser eifrige Liebhaber der Astronomie gerade zu dieser Zeit die Sonne wirklich beobachtet hatte und so durch wirkliche Erfahrung \SourceLeaf{627} bestätigen konnte, was man auch nach unseren sonstigen heutigen Kenntnissen gar nicht anders erwarten konnte, dass nämlich vom lockeren Kometenkörper, obwohl er sich gewiss und genau zwischen Erde und Sonne befand, auf der Sonnenscheibe nicht die allergeringste Spur erkannt wurde. Wir erinnern uns, in öffentlichen Blättern gelesen zu haben, dass noch ein anderer Beobachter in Wien, der in derselben Stunde ebenfalls die Sonne beobachtete, gar nichts Besonderes wahrgenommen hat. --- Beobachtungen dieses Kometen in der zweiten Augusthälfte von Dr. \textsc{Olbers} sowie Elemente seiner Bahn von Herrn \textsc{Bouvard} und von Herrn \textsc{Dirksen}. --- Über die Aufgabe, die Position eines Gestirns anhand beobachteter Alignements zu bestimmen, von Dr. \textsc{Olbers}. Im Jahrbuch für 1821 findet sich eine Lösung dieser Aufgabe von \textsc{Bessel}. In der von \textsc{Olbers} gewählten Form braucht man zwei Logarithmen weniger. Es ist jedoch zu bedenken, dass die dafür erforderlichen 20 Logarithmen an ebenso vielen verschiedenen Stellen stehen, während die für die \textsc{Bessel}'sche Form nötigen 22 nur an 17 verschiedenen Stellen nachzuschlagen sind. Wir bemerken außerdem, dass man bei zusätzlicher Verwendung der bekannten Hilfstafel für Logarithmen von Summen und Differenzen mit 18 Logarithmen an 16 Stellen oder mit 20 Logarithmen an 14 Stellen auskommt. --- Beobachtungen des Kometen vom Juli 1819, Beschreibung des sechs Fuß langen \textsc{Reichenbach}'schen Mittagsfernrohrs sowie Beobachtung der Jupiteropposition und des Polarsterns von Hofrat \textsc{Gauss}. Bei dem, was hier über die absolute Zuverlässigkeit der Zenitdistanzen gesagt wird, weisen wir auf zwei Druckfehler hin, die den Sinn ganz verfälschen. Oben auf S. 239 [S. 427] muss nämlich statt merklich \emph{unmerklich} stehen, und in der 5. Zeile von unten muss statt doch kann ich schon glauben \emph{doch kann ich schwer glauben} gelesen werden. --- Beobachtungen des Kometen vom Juli 1819 an der Krakauer Sternwarte von Prof. \textsc{Leski}. --- Verzeichnis der genauen Länge und Breite von neun der bedeutendsten Fixsterne, für die im Nautical Almanac und in der Connoissance des tems Mondabstände angegeben sind, entnommen aus dem Jahrgang 1821 der letzteren (siehe die Anzeigen des Vorjahres, S. 1907 [S. 615 dieses Bandes]). --- Beobachtungen eines neuen Kometen in Marseille und Elemente seiner Bahn, berechnet von Prof. \textsc{Enke}. Dieser Komet wurde, soweit bisher bekannt ist, nur in Marseille im Juni 1819 beobachtet, und es ist sehr bemerkenswert, dass sich die Beobachtungen durch eine Ellipse mit einer $2\frac{1}{2}$-jährigen
Umlaufzeit viel besser darstellen lassen als durch eine Parabel. --- Beobachtungen des (anderen) Kometen von 1819 an der Greenwicher Sternwarte und daraus von Herrn \textsc{Rumker} abgeleitete parabolische Elemente. --- Die kurzen Nachrichten, die zum Schluss noch folgen, enthalten weitere astronomische Beobachtungen und Notizen, die hier nicht einzeln angeführt werden können. Auf der beigefügten Kupfertafel ist der geozentrische Lauf des Kometen vom Juli 1819 und eine Projektion seiner wahren Bahn sowie die der Bahn des \textsc{Enke}'schen und des anderen von November 1818 bis Januar 1819 beobachteten Kometen dargestellt.
\par\clearpage\endgroup
